%%%%%%%%%%%%%%%%%%%%%%%%%%%%%%%%%%%%%%%%%%%%%%%%%%%%%% 
%%%%%%%%%%%%%%%%%%%%%%%%%%%%%%%%%%%%%%%%%%%%%%%%%%%%%%
%
% DEFINITIONS GENERALES DU DOCUMENT
%
%%%%%%%%%%%%%%%%%%%%%%%%%%%%%%%%%%%%%%%%%%%%%%%%%%%%%%
%%%%%%%%%%%%%%%%%%%%%%%%%%%%%%%%%%%%%%%%%%%%%%%%%%%%%%

\documentclass[a4paper,twoside,12pt]{report}           % En deux faces

%\pdfoutput=1 (dvi backward search) 1 (pdf c'est tout}
\pdfoutput=1

\usepackage{graphicx}

\ifnum\pdfoutput<1
\DeclareGraphicsExtensions{.eps}

\newcommand{\inputxfig}[1]{\includegraphics{#1.eps}}
\else
\DeclareGraphicsExtensions{.pdf}
\newcommand{\inputxfig}[1]{\input{#1.pdftex_t}}
\fi


%#{latexmk}include.sources.graphiques{{figures/}}
\graphicspath{{figure/}{./}}


\usepackage[T1]{fontenc}
\usepackage{amsthm,amssymb,amsmath,amsfonts}
\usepackage{latexsym}
\usepackage{ae}
\usepackage{aecompl}
\usepackage[cyr]{aeguill}

\usepackage[utf8]{inputenc}
\usepackage{color}

\def\stackunder#1#2{\mathrel{\mathop{#2}\limits_{#1}}}

%%%%%%%%%%%%%%%%%%%%%%%%%%%%%%%%%%%%%%%%%%%%%%%%%%%%%%
\newtheorem{theo}{Théorème}
\newcommand{\fmatlab}[1]{\underline{\textnormal{Fonctions Matlab utilisées~:}}
  \textnormal{\texttt{#1}}\vspace{0.3cm}}

\newcommand{\exintro}[0]{\upshape\rmfamily}
\newcommand{\exquestion}[0]{\itshape\sffamily}
\newcommand{\script}[0]{\normalfont\ttfamily\bigskip}




%%%%%%%%%%%%%%%%%%%%%%%%%%%%%%%%%%%%%%%%%%%%%%%%%%%%%%
%%%%%%%%%%%%%%%%%%%%%%%%%%%%%%%%%%%%%%%%%%%%%%%%%%%%%%
%
% MISE EN PAGE
%
%%%%%%%%%%%%%%%%%%%%%%%%%%%%%%%%%%%%%%%%%%%%%%%%%%%%%%
%%%%%%%%%%%%%%%%%%%%%%%%%%%%%%%%%%%%%%%%%%%%%%%%%%%%%%


%\setlength{\hoffset}{1cm}
%\setlength{\textwidth}{16cm}
\addtolength{\textwidth}{3cm}

%\setlength{\oddsidemargin}{0cm}
%\setlength{\evensidemargin}{0cm}
\addtolength\oddsidemargin{0cm}
\addtolength\evensidemargin{-3cm}
%\addtolength\marginparwidth{0.5cm}

\addtolength\topmargin{-1.5cm}
\addtolength\textheight{3cm}


%%%%%%%%%%%%%%%%%%%%%%%%%%%%%%%%%%%%%%%%%%%%%%%%%%%%%%
%%%%%%%%%%%%%%%%%%%%%%%%%%%%%%%%%%%%%%%%%%%%%%%%%%%%%%
%
% Extensions
%
%%%%%%%%%%%%%%%%%%%%%%%%%%%%%%%%%%%%%%%%%%%%%%%%%%%%%%
%%%%%%%%%%%%%%%%%%%%%%%%%%%%%%%%%%%%%%%%%%%%%%%%%%%%%%

\usepackage{graphicx}
\usepackage{lscape}
\usepackage{longtable}

\usepackage[french]{babel}             % distribution french
\usepackage[french]{varioref}



%%%%%%%%%%%%%%%%%%%%%%%%%%%%%%%%%%%%%%%%%%%%%%%%%%%%%%
%%%%%%%%%%%%%%%%%%%%%%%%%%%%%%%%%%%%%%%%%%%%%%%%%%%%%%
%
% Utilisation de fancyheadings.sty
%
%%%%%%%%%%%%%%%%%%%%%%%%%%%%%%%%%%%%%%%%%%%%%%%%%%%%%%
%%%%%%%%%%%%%%%%%%%%%%%%%%%%%%%%%%%%%%%%%%%%%%%%%%%%%%
\newfont{\scaledfont}{cmr12 scaled 4000}
\renewcommand{\thechapter}{\arabic{chapter}}
\makeatletter
\def\@makechapterhead#1{
  {\parindent \z@ \raggedright \normalfont
    \interlinepenalty\@M
    \parbox{0.70\textwidth}{\Huge \bfseries #1}\ %
    \hfill\scaledfont{TP\thechapter}\par\normalsize%
    \rule{\textwidth}{1pt}%  <--- the rule
    \nobreak
    \vskip 40\p@
  }%
}
\makeatother


\addto\captionsfrench{%
\renewcommand{\chaptername}{Séance}
%\renewcommand{\chaptermark}[1]{\markboth{\chaptername~\thechapter~:#1}{}} %
}

\usepackage{fancyhdr}
\pagestyle{fancy}


\addtolength{\headheight}{0.2\baselineskip} % Pour éviter le overfull \vbox .... has occured while ouput  is active voir doc fancyhdr

%\renewcommand{\chaptermark}[1]{\markboth{\chaptername~\thechapter~: \ #1}{}}
%\renewcommand{\sectionmark}[1]{\markright{ \thesection  \ #1}}

%\renewcommand{\chaptermark}[1]{\markboth{\textit{#1}}{}}
%\renewcommand{\sectionmark}[1]{\markright{\textit{\thesection \ #1}}}
%\renewcommand{\verb}[1]{\footnotesize {\ttfamily #1}}

%%%%%%%%%%%%%%%%%
%% PAGES DE GAUCHE (PAGES PAIRES) Cas de deux faces
%%%%%%%%%%%%%%
\fancyhead[LE]{\thepage}                                        % Indications en haut à gauche
\fancyhead[CE]{}                                                % Indications en haut au centre
\fancyhead[RE]{\leftmark}                                       % Indications en haut à droite
\fancyfoot[LE]{\tiny Analyse et traitement du signal - \today}     % Indications en bas à gauche
\fancyfoot[CE]{}                                                % Indications en bas au centre
\fancyfoot[RE]{}                                                % Indications en bas à doite
%%%%%%%%%%%%%%%%%
%% PAGES DE DROITE (PAGES IMPAIRES) Cas de deux faces
%%%%%%%%%%%%%%
\fancyhead[LO]{\rightmark}                                      % Indications en haut à gauche
\fancyhead[CO]{}                                                % Indications en haut au centre
\fancyhead[RO]{\thepage}                                        % Indications en haut à droite
\fancyfoot[LO]{}                                                % Indications en bas à gauche
\fancyfoot[CO]{}                                                % Indications en bas au centre
\fancyfoot[RO]{\tiny Analyse et traitement du signal - \today}     % Indications en bas à doite


%%%%%%%%%%%%%%%%%
%% PAGES DE GAUCHE (PAGES PAIRES) Cas de une face
%%%%%%%%%%%%%%
%\fancyhead[LE]{}    % Indications en haut à gauche
%\fancyhead[CE]{}   % Indications en haut au centre
%\fancyhead[RE]{}    % Indications en haut à droite
%\fancyfoot[LE]{}     % Indications en bas à gauche
%\fancyfoot[CE]{}    % Indications en bas au centre
%\fancyfoot[RE]{}      % Indications en bas à doite
%%%%%%%%%%%%%%%%%
%% PAGES DE DROITE (PAGES IMPAIRES) Cas de une face
%%%%%%%%%%%%%%
%\fancyhead[LO]{\rightmark}    % Indications en haut à gauche
%\fancyhead[CO]{}   % Indications en haut au centre
%\fancyhead[RO]{\thepage}    % Indications en haut à droite
%\fancyfoot[LO]{\leftmark}     % Indications en bas à gauche
%\fancyfoot[CO]{}    % Indications en bas au centre
%\fancyfoot[RO]{\tiny Sébastien \textsc{Bosch} - LAAS-RIA - \today}      % Indications en bas à doite

\renewcommand{\headrulewidth}{0.4pt}
%\renewcommand{\footrulewidth}{0.4pt}




%\input{tcilatex}

\begin{document}


% \thispagestyle{empty}

% \begin{center}

%   \begin{tabular}[l]{ll}
%     \includegraphics[width=0.2\textwidth]{Figure2}\hspace{0.5cm}

%     &\vspace{-1.3cm} \\

%     &\textsc{Ministère de l'Enseignement Supérieur et de la
%       Recherche}\\
%     &\vspace{0cm}\textsc{Département des Sciences et
%       Techniques pour l'Ingénieur}

%   \end{tabular}

% \end{center}

% \vspace{-2cm}

% \begin{center}
%   {\LARGE \textsc{Travaux Pratiques~:}}

%   {\Large Filtrage numérique}


%   \vspace{2cm}


%   \includegraphics[width=0.8\textwidth]{Figure1}


%   \vspace{2cm}

%   {\large \textit{Grand Cru 2026--2027}}
% \end{center}

% \begin{tabular}[c]{lll}
%   \textbf{Sommeliers(ère)~:} \hspace{2cm}& Pascal \textsc{Acco} \\ 
%   \\
% \end{tabular}
% \pagebreak

% \begin{center} {\LARGE \textbf{Avant-propos}}
% \end{center}

% \sffamily

% Les TPs ne sont pas évalués directement, ils servent à appuyer la
% compréhension.

% Pour maximiser l'impact réel de ces enseignements, vous devez être
% complètement déconnectée (téléphone, montre, etc.) et effectuer un
% rapport manuscrit, car l'association d'une activité de motricité
% fine comme l'écriture manuelle renforce la mémorisation.


% \subsection*{Contenu du TP}

% \subsection*{Déroulement du TP}


% \rmfamily
%% ___________________________________________
\chapter{Échantillonnage de signaux}
%% __________________________________________


Nous allons travailler de trois manières~:
\begin{description}
\item Expérimentalement -- en utilisant un GBF, un microcontrôleur et
  un oscilloscope.
\item Graphiquement -- en utilisant un papier millimétré, un compas,
  une règle et une échelle logarithmique
\item Analytiquement -- lorsque les solutions analytiques existent, on
  les exprimera. L'application numérique est faite graphiquement, et
  on confrontera avec les résultats expérimentaux.
\end{description}


\section{Échantillonnage d'un sinus}
\exintro{}
\label{part:echantillonnage}

Le micro-contrôleur possède~:
\begin{itemize}
\item une entrée analogique, notée $x(t)$, et est déjà programmé pour
  récupérer les échantillons mesurés par un Convertisseur Analogique
  Numérique (CAN ou ADC en anglais) avec une période $T_e$ fixe. Le
  processeur stocke périodiquement ces échantillons, notés $x[k]$,
  pour effectuer un ou des traitements dits \emph{numériques}.
\item Une sortie de pontage (\emph{bypass} en anglais), notée $a[k]$.
  La sortie de pontage est effectuée le plus rapidement possible, elle
  consiste à recopier la valeur $x[k]$ du CNA directement vers un
  Convertisseur Numérique Analogique (CNA ou DAC en anglais). Ce
  traitement est très proche d'un gain unitaire~: $a[k] = 1 . x[k]$
  donc $G(z) = 1$. Le CNA offre une sortie à temps continu en
  maintenant constante la valeur de sortie, notée $a(t)$. Cette
  opération est appelée \emph{Bloqueur d'Ordre Zéro} (BOZ ou ZOH en
  anglais). La sortie $a(t)$ du BOZ du CNA est un signal en forme
  d'escalier que l'on peut observer à l'oscilloscope en branchant une
  sonde (en grip fil) sur la broche de rangée 1 et colonne 1 de la
  matrice de câblage de la carte électronique du microcontrôleur.
\item Une sortie filtrée, notée $b[k]$. Le processeur stocke dans un
  tampon d'entrée (buffer en anglais) les derniers échantillons
  $x[k]$, et a stocké dans un tampon de sortie les derniers
  échantillons $b[k]$. Il effectue le plus rapidement possible le
  calcul d'une récurrence (moyenne pondérée par des coefficients de
  ces deux tampons) pour obtenir l'échantillon de sortie $b(k)$. On a
  alors une fonction de transfert discrète inconnue
  $b(z)=G(z).x(z)=G(z).a(z)$. Cette sortie discrète $b[k]$ est envoyée
  sur une seconde sortie du CNA. Le signal continu $b(t)$ est
  observable à l'oscilloscope en branchant une sonde sur la rangée 1,
  colonne 2 de la matrice de câblage.
\end{itemize}

\textbf{Commencez votre rapport manuscrit en donnant le schéma du
  système de traitement.} Vous y indiquerez le signal entrant, le CAN
, le processeur qui calcule les deux signaux discret de sortie et les
deux CNA. Sur le schéma vous reporterez les signaux $x(t)$, $x[k]$,
$a[k]$, $a(t)$, $b[k]$, $b(t)$ et laisserez la place pour insérer,
plus tard, l'allure d'un chronogramme et l'expression analytique pour
chaque signal.

\subsection{Préliminaire GBF dans $0-3,3V$}
Il ne faut pas griller la carte et prendre le réflexe de toujours
vérifier les appareils et les réglages d'entrée avant de poursuivre le
montage.

Pour cela, branchez directement la sortie du GBF sur l'oscilloscope
(avec un BNC) et effectuer les réglages pour être sûre que la tension
de sortie du GBF est entre 0 et 3,3V.

\textbf{ La plage de tension $0-3.3V$ doit être respectée au risque
  d'endommager le CNA du $\mu$-contrôleur~!}

\begin{enumerate}
\item Essayez de produire et observer une sinusoïde à $100 Hz$ dont
  vous noterez l'offset et l'amplitude choisie (pas trop petite au
  moins 0,5 V) vous permettant de respecter la plage 0-3.3V.
\item Utilisez un second câble BNC pour relier la sortie SYNC/EXT du
  GBF\footnote{Certains GBF ont cette sortie désactivée. Il faut aller
    dans le menu \emph{utility -> synchronisation} pour activer la
    sortie} à l'entrée de déclenchement externe de l'oscillo. (EXT
  pour \emph{external trigger} en anglais) de manière à avoir une
  synchronisation sur un sinus (amplitude nulle à l'instant t=0)
  indépendante de l'offset.
\end{enumerate}

\textbf{Une fois ces réglages faits, il ne faut plus toucher à
  l'amplitude ni au décalage (offset) !}

\subsection{Échantillonnage d'un sinus lent à 100Hz}

On veut compléter notre schéma précédant en donnant les formes d'ondes
observées, mesurer la période d'échantillonnage $T_e$, et identifier
et mesurer les retards entre les signaux.

\subsubsection{Expérimentalement}

Le matériel est distribué sous forme de pochettes et vous disposez
d'un schéma de montage dans un classeur.
\begin{enumerate}
\item Effectuez le câblage (schéma et explication en TP) : Alimentez
  la carte $\mu$-contrôleur via son câble USB que vous brancherez sur
  un port USB quelconque (celui de l'oscilloscope ou du GBF ou d'un
  PC)~; connectez son entrée analogique avec un cordon Jack-BNC (à
  gauche de la carte) au sinus du GBF~; connectez la sonde en voie 2
  sur la sortie de pontage (rangée 1 colonne 1 de la matrice). Pour
  cela, il faudra construire le cordon jack-BNC avec un câble
  Jack-bananes et un plot BNC-Bananes.
\item observez $a(t)$ et $x(t)$ et en déduire la période
  d'échantillonnage $T_e$~;
\item modifiez légèrement la fréquence du BGF pour stabiliser $a(t)$,
  essayez de comprendre le phénomène (indice~: aucun oscillateur n'est
  parfait, et il y a deux oscillateurs dans ce montage)~;
\item une fois $a(t)$ stabilisé (ou juste en appuyant sur le STOP de
  l'oscilloscope). Représentez graphiquement ces deux signaux observés
  et mesurez le retard entre $a(t)$ et $x(t)$ observez. Identifiez
  leurS causeS possibleS.
\item indiquez sur ce schéma comment mesurer les échantillons $a[k]$
  en observant $a(t)$ en indiquant sous forme de bâtonnets le signal
  discret $a[k]$ par-dessus les signaux continus $a(t)$ et $x(t)$
\end{enumerate}


\subsubsection{Analytiquement}

L'échantillonnage d'un signal continu $t\mapsto s(t)$ (où $s$ est une
fonction) consiste à passer au signal discret $k \mapsto s[k]$ (où $s$
est une suite numérique) en considérant que l'instant $t$ d'un
échantillon $k$ est $t=k.T_e$.
\begin{enumerate}
\item \'Ecrivez la forme analytique de $x(t)$ pour notre sinus
  synchronisé avec une amplitude $A$ et un décalage noté $x_0$, où
  vous ferez apparaitre la fréquence $f$.
\item Sachant que la période de la fonction sinus est de $2\pi$ et que
  son argument est sans dimension. Retrouvez le lien entre la
  fréquence $f$, la pulsation $\omega$ et la période $T$ du sinus.
\item Donnez l'expression de $a[k]=x[k]$ à partir de $x(t)$. Comme $k$
  est un numéro d'échantillon, il est sans dimension. Vérifiez que
  l'argument de votre fonction sinus est toujours sans dimension. Vous
  pouvez donner cette expression en fonction de $f$ et $F_e$~: en
  discret le ratio $\frac{f}{F_e}$ est appelé la fréquence normalisée
  et est essentielle.
\item On peut observer exactement le même signal discret $a[k]$ pour
  d'autres fréquences $f$ et $F_e$. Donnez tous les couples $(f,F_e)$
  permettant d'obtenir ce signal. Conclure.
\end{enumerate}

On va représenter le sinus avec un peu plus d'imagination~: la
représentation par phaseur, ou amplitude complexe~:

\begin{equation}
  \label{eq:phaseur}
  s(t) = \mathcal{R}\left\{\hat{z}.e^{i.\omega.t}\right\}
\end{equation}

Par exemple, si on prend le phaseur $\hat{z_i}=1 = 1.e^{i.0}$ on
obtient le signal cosinus~:
\begin{equation}
  \label{eq:Zi}
  s_i(t) = \mathcal{R}\left\{z_i.e^{i.\omega.t}\right\} = \mathcal{R}\left\{e^{i.\omega.t}\right\}=\cos(\omega.t)
\end{equation}


On parle de signal \og{} In phase\fg{}, car il est en phase avec le
cosinus pris pour référence, d'où le nom du phaseur $\hat{z_i}$.

\begin{enumerate}
\item Donnez l'expression réelle d'un signal de phaseur quelconque
  $\hat{z}=\rho.e^{i.\varphi_0}$ et trouvez ainsi le lien entre
  l'amplitude et la phase du signal et le phaseur~;
\item Quelle amplitude et quelle phase faut-il apporter au cosinus
  pour obtenir un sinus ?  En déduire le phaseur $\hat{z_q}$ d'un
  sinus d'amplitude 1, et vérifiez que cela fonction avec
  l'eq.\ref{eq:phaseur}.
\item Donnez le phaseur $z_{100}$ qui représente votre signal
  sinusoïdal à 100 Hz d'amplitude choisie. Comment représenter son
  décalage ?
\end{enumerate}

\subsubsection{Graphiquement}

Si l'on représente $e^{i.\omega.t}$ dans le plan complexe, on obtient
le cercle unité avec pour point de départ la valeur 1 obtenue pour
$t=0$.

Si l'on représente $\hat{z}.e^{i.\omega.t}$ dans ce plan, on obtient
le cercle de rayon $\rho$ avec un point de départ qui vaut $\hat{z}$
obtenu pour $t=0$.

Prenons notre signal à la fréquence $f_{100}=100$Hz, de pulsation
$\omega_{100}$ et de période $T_{100}$

\begin{enumerate}
\item Représentez, graphiquement dans le plan complexe, votre signal
  $\hat{z_{100}}.e^{i.\omega_{100}.t}$ avec son point de départ, sa
  valeur pour $t=\frac{T_{100}}{2}$, $t=\frac{T_{100}}{4}$~;
\item Donnez les numéros d'échantillons $k$ correspondant à
  $t=0$,$t=\frac{T_{100}}{4}$, $t=\frac{T_{100}}{2}$ et $t+T$~~: en
  déduire la période de la suite numérique $\left\{a[k]\right\}$
\end{enumerate}

\subsection{Sous-échantillonnage = périodicité du spectre}

On vient de voir que l'onde pure $x(t)$ à 100 Hz échantillonnée à 10
kHZ donne un signal numérique $a[k]$ de période 100.  Cela est vrai
pour tout échantillonnage de fréquence $F_e$ d'une onde de fréquence
$f$.

Nous allons voir qu'une infinité d'autres fréquences normalisées
peuvent donner ce même signal~: le sous-échantillonnage.

\subsubsection{Expérimental}
\begin{enumerate}
\item Ne pouvant pas changer $T_e$ (c'est programmé dans le
  $\mu$-contrôleur), utiliser vos équations de $a[k]$ précédentes pour
  trouver la fréquence, notée $f_6$, du GBF qui fait que le signal
  échantillonné soit 6-périodique.  Sans rien changer au montage,
  régler votre GBF pour observer ce signal, noté $a_6(t)$, stable
  (synchroniser par tâtonnement sur la fréquence du GBF). Mesurez
  (rapidement avec les carreaux de l'oscilloscope) et donnez le signal
  6-périodique $a_6$ sous la forme
  $(\!( a[0], a[1], a[2], \dots, a[5])\!)$ où les double parenthèse
  indiquent une suite périodique.
\end{enumerate}

Maintenant essayons des fréquences plus rapides~:
\begin{enumerate}
\item augmentez progressivement la fréquence $f$ du signal GBF en
  passant par $\frac{F_e}{2}$, $F_e$ puis la fréquence $f=f_6 + F_e$~;
\item représentez le signal temporel $x(t)$ et sa discrétisation
  $a[k]$ (petits ronds des bâtonnets par-dessus la courbe) pour la
  fréquence $\widetilde{f}=\frac{1}{6}$, puis pour la fréquence
  $\widetilde{f}=\frac{1}{6} + 1$, puis plus loin~: vous observez le
  phénomène de sous-échantillonnage.
\item pour un fort sous-échantillonnage (10 fois $F_e$)
\end{enumerate}

Vous venez de constater que tous signaux continus de fréquences égales
modulo $F_e$ donnent exactement le même signal discret.

\subsubsection{Analytiquement}
Vous avez vu que la fréquence $f_6$ telle que
$\widetilde{f_6}=\frac{1}{6}$ donne un signal discret 6-périodique.

Supposez qu'il existe une autre fréquence que $f_6$ donnant le même
signal, que l'on modélise par $f = f_6 + \Delta_f$. Cela correspond à
un écart de fréquence normalisé noté
$\widetilde{f} = \widetilde{f_6} + \widetilde{\Delta_f}$. En déduire
la contrainte sur $\widetilde{\Delta_f}$ pour que le signal ne change
pas.
\begin{enumerate}
\item Donnez toutes les fréquences normalisées équivalentes à $f_6$~;
\item Généralisez pour toute fréquence $f$
\end{enumerate}

On verra que la Transformée de Fourier d'un Signal Discret (TFSD)
représente un signal discret par une somme de signaux
$w_f=e^{i.2\pi.f.k.T_e}$ pour des fréquences $f\in\mathbb{R}$.

Comme les fréquences $f+m.F_e$ sont toutes équivalentes et donnent le
même signal $w_f$, alors les composantes de la TFSD sont toutes les
mêmes pour ces fréquences~:

\textbf{Un spectre discret (TFSD) est forcément $F_e$ périodique.}

\subsection{Fréquences négatives = symétrie du spectre}

Une autre représentation des ondes pures par des complexes, utilise
les fréquences négatives. En effet, une onde à une fréquence négative
$-f$ est le conjugué de l'onde positive~:
$e^{i.2\pi.(-f).t}=e^{-i.2\pi.f.t}=\overline{e^{i.2\pi.f.t}}$.

Ainsi, la partie réelle utilisée par un phaseur peut s'obtenir par
addition de deux conjuguées~:
\begin{equation}
  \label{eq:cn}
  \mathcal{R}\left\{\hat{z}.e^{i.\omega.t}\right\}=\frac{\hat{z}.e^{i.\omega.t}+\overline{\hat{z}.e^{i.\omega.t}}}{2} = \underbrace{\frac{\hat{z}}{2}}_{c(\omega)}.e^{i.\omega.t} + \underbrace{\frac{\overline{\hat{z}}}{2}}_{c(-\omega)}.e^{i.(-\omega).t}
\end{equation}

Il s'agit bien de la représentation en coefficient $c(\omega)$
utilisée dans les transformées et séries de Fourier.

\textbf{Ainsi tout spectre d'un signal réel à une symétrie de
  Hilbert~: $c(-f)=\overline{c(f)}$.}

Les fréquences négatives d'ondes complexes sont utilisée pour
représenter des ondes réelles.

\subsubsection{Graphiquement}

On va représenter le signal complexe (fonction de $\mathbb{R}$ dans
$\mathbb{C}$)~:
 
$s_a[k]=e^{i.2\pi.f.\underbrace{k.Te}_{t}}=e^{i.2\pi.\underbrace{\frac{f}{F_e}}_{\widetilde{f}}.k}$

Ce signal est appelé \emph{signal analytique} car il permet comme un
phaseur de retourner au signal réel avec
$a[k]=\mathcal{R}\left\{s_a[k]\right\}$

Nous allons représenter ce signal analytique dans le plan complexe
pour $\widetilde{f}=\frac{1}{6}$.
   
\begin{enumerate}
\item Sur un papier millimétré, représentez dans le plan complexe les
  7 premiers échantillons de $s_a[k]=A.e^{i.\frac{2\pi}{6}.k}$ (pour
  construire l'angle $\frac{\pi}{3}$ au compas, il suffit de
  construire un triangle équilatéral~: la rosace dans le cercle de
  l'école primaire...).  Où $A$ est l'amplitude choisie pour votre
  signal.
\item on note $q=e^{i\frac{2\pi}{6}}$, mesurez la partie réelle et la
  partie imaginaire de $q$ et retrouvez analytiquement ce que sont ces
  valeurs (c'est un triangle équilatéral dont vous cherchez la hauteur
  et la demi-base...).  Exprimez le signal analytique 6-périodique
  $s_a[k] = (\!( 1, ..., )\!)$ en fonction de $q$. On doit trouver une
  suite géométrique.
\item Mesurez les parties réelles et retrouvez les mesures
  expérimentales $\left\{a_6[k]\right\}=(\!(0,a[1], \dots, )\!)$
  faites à l'oscillo.
\item Refaites de même pour la fréquence négative
  $\frac{-1}{6}$. Remarquez que l'on a le conjugué de la séquence
  analytique précédente qui nous fait tourner à l'envers dans le
  cercle. Remarquez temporelle notée $a_{-6}$ est la même que celle de
  $a_6$, mais jouée à l'envers~!
\end{enumerate}

On obtient ainsi que pour les fréquences négatives, le signal
analytique (et donc la TFSD) est le conjugué de la fréquence
positive. C'est la symétrie de Hilbert.

\textbf{Pour les fréquences négatives d'une onde pure $w_{-f}(t)$ (ou
  en discret $w_{-f}[k]$), on obtient le signal antipodal
  $w_{-f}(t) = w_f(-t)$ (ou $w_{-f}[k]=w_f[-k]$)} Pour ce qui est du
signal réel cela revient à obtenir le signal antipodal.

\subsubsection{Expérimentalement}


Difficile de demander à un GBF une fréquence négative
$\widetilde{f}=\frac{-1}{6}$. En revanche, on peut utiliser la
périodicité en $F_e$ du spectre pour trouver une fréquence équivalente
(donc exactement le même signal discret) qui soit positive.

\begin{itemize}
\item Placez le GBF sur une fréquence image de
  $\widetilde{f}=\frac{-1}{6}$ qui soit positive, stabilisez le signal
  $a(t)$ et stoppez-le lorsqu'il est en phase ($a[0]=0$)
\item mesurez rapidement (avec les carreaux de l'oscilloscope) le
  signal 6-périodique $a_{-6}=(\!( a_{-6}[0], \dots a_{-6}[5])\!)$
\item vous devriez vérifier $a_{-6}[k]=a_{6}[-k]=a_{6}[6-k]$~: le
  signal périodique joué à l'envers dans le temps.
\end{itemize}

\section{Spectre d'un signal discret}
 
Dans une TFSD, on décompose un signal discret (potentiellement avec
une part imaginaire) en somme de signaux $w_f[k]=e^{i.2\pi.f.k.T_e}$
pour tout $f$ dans $\left[0,\; F_e \right[$ uniquement ! Car au-delà,
chaque signal est périodiquement équivalent à son signal modulo
$F_e$~: $w_{f+M.F_e}=w_f$.

Dans le cas de signaux réels, déterminer la TFSD du signal discret
$ k \mapsto s[k]$, c'est calculer chaque phaseur $\hat{z}(f)$ pour
chaque fréquence qui compose le signal~:
\begin{equation}
  \label{eq:tfsd_decompo}
  s[k] = \int\limits_{0\leq f < F_e}{\mathcal{R}\left\{\widehat{z_s}(f).\underbrace{e^{i.2\pi.f.k.T_e}}_{w_f[k]}\right\}.df}
\end{equation}

\subsection{Axe des fréquences}

\begin{itemize}
\item Dessinez sur un axe gradué l'axe des fréquences $f$ de $-F_e$ à
  $2.F_e$~;
\item Dessinez juste dessous l'axe des fréquences normalisées
  $\widetilde{f}$~;
\item Reportez encore dessous les signaux $w_0$, $w_2$, $w_6$,
  $w_{100}$ respectivement $1$, $2$, $6$ et $100$ périodiques devant
  les bonnes fréquences~;
\item Reportez les signaux antipodaux (fréquences négatives) de $w_0$,
  $w_2$, $w_6$, $w_{100}$ notés $w_{-0}$, $w_{-2}$, $w_{-6}$ et
  $w_{-100}$ devant les bonnes fréquences~;
\item Reportez aussi les signaux conjugués de type
  $\overline{w_\bullet}$ et les conjugués antipodaux.
\end{itemize}
 

 \subsection{Fréquence absolument la plus rapide = Nyquist}
 On voit que lorsque la fréquence $f$ augmente en continu, l'onde
 discrète n'est pas forcément plus rapide, par exemple un signal de
 fréquence $F_e$ est en fait un signal constant donc le plus lent
 possible.

 On se pose alors la question de la plus haute fréquence que l'on
 puisse avoir sur un signal discret.

 Il est utile de mettre le signal $e^{i.2\pi.f.k.T_e}$ sous la forme
 $e^{i.2\pi.\frac{f}{F_e}.k}=e^{i.\underbrace{2\pi.\widetilde{f}}_{\Delta_f}.k}=e^{i.\Delta_f.k}=q^k$
 avec $q=e^{i.\Delta_f}$ et où $Delta_f=2\pi\frac{f}{F_e}=\omega.T_e$
 est nommé l'incrément angulaire.

 \begin{enumerate}
 \item À quelles fréquences correspondent les incréments
   $\Delta=\pm\frac{\pi}{2}$, $\Delta=\pm\pi$ et
   $\Delta=\pm\frac{\pi}{2}$ ?
 \item montrez que l'incrément $\Delta=\pi + \epsilon$ (avec
   $\epsilon<\pi$) est en fait le signal antipodal de
   $\Delta=\pi - \epsilon$ qui est plus petit donc plus lent. En
   déduire l'incrément le plus rapide et la fréquence la plus rapide
   en discret~: la fréquence de Nyquist.
 \item Donnez la séquence discrète correspondant au signal le plus
   rapide et en déduire sa période
 \item Donnez une fréquence haute qui donne un signal qui tourne
   lentement, mais à l'envers, puis à l'arrêt, puis lentement en
   avant...
 \end{enumerate}

 Il s'agit du phénomène stroboscopique observé lorsque l'on observe
 les roues d'une voiture à la télévision, tournant vite puis lentement
 à l'envers, puis à l'arrêt puis lente en avant, etc.
 

 \subsection{Tracé de spectre}
 
 \subsubsection{Spectre de phaseurs}
 
 Le signal 6 périodique du TP de forme $x(t) = x_0 + A \sin(2\pi.f.t)$
 en continu et sa forme échantillonnée sont faciles à décomposer
 puisque l'on a déjà identifié les phaseurs de ce sinus.

 On peut tracer un graphe du module $\rho(f)$ du phaseur (sous forme
 de bâtonnets) en utilisant notre axe de fréquence. On ne fera pas le
 graphe de la phase ici.
 \begin{itemize}
 \item Identifiez leS fréquenceS des composantes de votre signal, et
   donnez le phaseur de chacune de ces composantes.
 \item Reportez sous forme d'un bâtonnet d'amplitude
   $\rho(f)=|\hat{z_s}(f)|$ sur toutes les fréquences où se trouve le
   signal composante discret $w_f$ et indiquez à côté la valeur
   complexe du phaseur.

 \item Normalement les deux composantes fréquentielles doivent donner
   7 bâtonnets entre $-F_e$ et $2.F_e$.
 \item Faites apparaitre la périodicité de votre spectre.
 \end{itemize}


 \subsubsection{Spectre complexe}
 On peut introduire les fréquences négatives pour décomposer le signal
 réel
 $\mathcal{R}\left\{\hat{z_s(f).w(f)}\right\}=
 \frac{z_s(f).w_f}{2}+\frac{z_s(-f).w_{-f}}{2}$ en somme de conjugués
 et avoir la décomposition en coefficients complexes
 $c(f)=\frac{\widehat{Z_s}(f)}{2}$ ~:
 \begin{equation}
   \label{eq:itfsd_cn}
   s[k] = \int\limits_{-F_e < f < F_e}{c(f).e^{i.2\pi.f.k.T_e}.df}
 \end{equation}
 
 \begin{itemize}
 \item Dessinez le spectre des modules $|c(n)|$. et indiquez à côté
   des bâtonnets la valeur complexe de $c(n)$. Il est d'usage de
   dessiner un bâtonnet imaginaire pur négatif comme un bâtonnet
   négatif, mais c'est un abus.
 \item Vous devez avoir maintenant 10 bâtonnets et pouvez faire
   apparaitre des axes de symétrie de Hilbert
 \item Que ce passe-t-il, si les coefficients ne respectent pas la
   symétrie de Hilbert ? On en conclu que la représentation en $c(f)$
   est plus générale que cette en phaseur limitée aux signaux réels.
 \end{itemize}

 
 \subsubsection{Spectre polaire}

 Si l'on note la forme polaire des phaseurs
 $\widehat{z_s}(f)=\rho(f).e^{i.\varphi(f)}$ dépendant de $f$, on
 obtient alors la recomposition (TFSD inverse) sous sa forme polaire~:
 \begin{equation}
   \label{eq:itfsd_polaire}
   s[k] = \int\limits_{0\leq f < F_e}{\rho(f).\cos\left(2\pi.f.k.T_e+\varphi(f)\right).df}
 \end{equation}

 Le spectre serait alors la courbe $\rho(f)$ (associée à l'amplitude)
 et la courbe $\varphi(f)$ (associée à la phase).

 \begin{itemize}
 \item Remarquez que le spectre tracé précédemment correspond au
   spectre d'amplitude.
 \item Ajoutez une deuxième courbe avec le spectre de phase
   $\varphi(f)$ sur la même échelle.
 \end{itemize}


 \subsubsection{Spectre en composante I et Q}

 Puisque nous avons le phaseur de cosinus (In phase) qui est réel pur
 $\hat{z_I}=1$ et celui du sinus (Quadrature phase) qui est imaginaire
 pur $\hat{z_Q}=-i$, on peut toujours décomposer
 $\hat{z_s}= a(f) .\hat{z_I} + b(f) .\hat{z_Q} = a(f) - i. b(f)$ pour
 avoir la représentation I/Q de Fourier~:
 \begin{equation}
   \label{eq:itfsd_polaire}
   s[k] = \int\limits_{0\leq f < F_e}{a(f).\cos\left(2\pi.f.k.T_e\right) + b(f).\cos\left(2\pi.f.k.T_e\right).df}
 \end{equation}
 Le spectre est alors une fonction $a(f)$ (associée aux composantes
 paires) et une fonction $b(f)$ (les impaires).
 \begin{itemize}
 \item Votre signal $x(t)$ est-il pair, impair ou quelconque ? En
   déduire si les $a(f)$ sont tous nuls, ou les $b(f)$ sont tous nuls,
   ou pas.
 \item Esquissez rapidement le spectre en bâtonnet $a(f)$ et $b(f)$
 \end{itemize}


\end{document}
